% Source: https://artofproblemsolving.com/community/c3h1642066_a_new_way_of_defining_lcms_maybe
% Stable identity: aops-topic-1642066-q01
The conjecture is false. Take $x=1$, $y=6$:
\[
1+\dfrac{1}{2}+\dfrac{1}{3}+\dfrac{1}{4}+\dfrac{1}{5}+\dfrac{1}{6}=\dfrac{60+30+20+15+12+10}{60}=\dfrac{147}{60}=\dfrac{49}{20},
\]
whose denominator $20$ is not $\operatorname{lcm}(1,2,\dots,6)=60$. (Another counterexample: $\dfrac{1}{3}+\dfrac{1}{4}+\dfrac{1}{5}+\dfrac{1}{6}=\dfrac{57}{60}=\dfrac{19}{20}$.)

What is true is weaker: if all the fractions are written over the least common denominator, that denominator is the lcm (by definition), so the reduced denominator always \emph{divides} $\operatorname{lcm}(x,\dots,y)$. Equality fails whenever the numerator shares a factor with the lcm, as $147=3\cdot49$ does with $60$. So the sum gives no new way to compute the lcm.
